\documentclass[11pt]{article}
\usepackage[margin=1in]{geometry}
\usepackage{amsmath,amssymb,amsthm,mathtools}
\usepackage{graphicx}
\usepackage{microtype}
\usepackage[hidelinks]{hyperref}

\newcommand{\R}{\mathbb R}
\newcommand{\B}{\mathbb B}
\newcommand{\dist}{\operatorname{dist}}
\newcommand{\Int}{\operatorname{Int}}
\newtheorem{theorem}{Theorem}
\newtheorem{proposition}{Proposition}
\newtheorem{lemma}{Lemma}
\newtheorem{remark}{Remark}

\title{Exact stability of affine sections of Hilbert balls\\
and the sharp role of uniform continuity}
\author{Candidate full extension for arXiv:1902.02363}
\date{19 July 2026}

\begin{document}
\maketitle

\begin{center}
\fbox{\parbox{0.91\textwidth}{
\textbf{Status: candidate full answer, likely valid.}
For closed balls in arbitrary Hilbert spaces, affine fibers of a closed-range
operator have an explicit orthogonal-product form and an exact Hausdorff
distance.  This removes compactness and ambient finite dimensionality from the
source example.  Arbitrary continuous objectives remain valid when the kernel
is finite-dimensional; uniformly continuous objectives work for every kernel.
An explicit bounded continuous objective shows that mere continuity is not
enough when the kernel is infinite-dimensional.}}
\end{center}

\section{Source question}

The source is Daniel Reem, Simeon Reich, and Alvaro De Pierro,
\emph{Stability of the optimal values under small perturbations of the
constraint set}, arXiv:1902.02363v4 \cite{RRD}.  Example 5.11 proves continuity
of
\[
 t\longmapsto \inf\{f(x):x\in C,\ Lx=t\}
\]
when $C\subset\R^n$ is compact, convex, and has nonempty interior and
$f:C\to\R$ is continuous.  In Remark 5.12(i), page 19, the authors write:

\begin{quote}
``We believe that the theory developed in this section can be extended
further, and, in particular, that it is possible to remove (at least in some
interesting cases) the compactness and finite-dimensionality assumptions from
Example 5.11.''
\end{quote}

\begin{figure}[ht]
\centering
\includegraphics[width=\textwidth]{figures/open_problem_crop.png}
\caption{The source question in full-width context, page 19.}
\label{fig:question}
\end{figure}

The result below supplies such a case, with an exact geometric modulus, and
also determines precisely where the source proof's use of compactness cannot
simply be discarded.

\section{Proof intuition}

Let $N=\ker L$ and $M=N^\perp$.  Closed range makes the restriction
$L|_M$ continuously invertible onto $L(H)$.  Thus fixing $Lx=t$ fixes
the $M$-coordinate of $x-c$, while the $N$-coordinate ranges over a
ball.  Pythagoras gives both its center and radius.  Two fibers are therefore
orthogonal translates of concentric kernel-balls, so their Hausdorff distance
is the Euclidean combination of center displacement and radius displacement.

This exact formula yields continuity, local Lipschitz bounds, and value
stability for uniformly continuous objectives without any compactness.  If
$N$ is finite-dimensional, the fibers encountered along a convergent
parameter sequence lie in one compact set, so mere continuity is again
enough.  If $N$ is infinite-dimensional, orthogonal directions allow
localized continuous wells to escape without accumulating.  Those wells
make the value function jump even though the fibers converge in Hausdorff
distance.

\section{Exact fiber geometry}

Let $H,Y$ be real Hilbert spaces and let $L:H\to Y$ be bounded and linear
with closed range $R=L(H)$.  Put
\[
 N:=\ker L,\qquad M:=N^\perp,
\]
and let
\[
 G:R\longrightarrow M
\]
be the inverse of $L|_M:M\to R$.  This inverse is bounded by the bounded
inverse theorem \cite{Conway}.  Fix $c\in H$, $r>0$, and the closed ball
\[
 C:=\overline B(c,r).
\]
For
\[
 J:=\{t\in R:\|G(t-Lc)\|\le r\},\qquad
 I:=\{t\in R:\|G(t-Lc)\|<r\},
\]
define
\[
 A_t:=C\cap L^{-1}(t),\quad
 z_t:=G(t-Lc),\quad
 \rho_t:=\sqrt{r^2-\|z_t\|^2}\qquad(t\in J).
\]
Here $I$ is exactly the set of parameters whose affine fiber meets
$\Int C$.

\begin{theorem}[Exact affine-section formula]\label{thm:geometry}
For every $t\in J$,
\begin{equation}\label{eq:fiber}
 A_t=c+z_t+\rho_t\B_N,
 \qquad \B_N:=\{n\in N:\|n\|\le1\}.
\end{equation}
If $N\ne\{0\}$, then for all $s,t\in J$,
\begin{equation}\label{eq:exactH}
 D_H(A_s,A_t)^2
 =\|z_s-z_t\|^2+|\rho_s-\rho_t|^2.
\end{equation}
If $N=\{0\}$, then $A_t=\{c+z_t\}$ and
\begin{equation}\label{eq:trivialH}
 D_H(A_s,A_t)=\|z_s-z_t\|.
\end{equation}
Consequently $t\mapsto A_t$ is Hausdorff-continuous on $J$ and locally
Lipschitz on $I$.  Quantitatively, for every $0\le b<r$ and all $s,t$
such that $\|z_s\|,\|z_t\|\le b$,
\begin{equation}\label{eq:localLip}
 D_H(A_s,A_t)
 \le \frac{r\|G\|}{\sqrt{r^2-b^2}}\,\|s-t\|.
\end{equation}
On all of $J$, including its boundary,
\begin{equation}\label{eq:holder}
 D_H(A_s,A_t)
 \le 2\sqrt{r\|G\|\,\|s-t\|}.
\end{equation}
\end{theorem}

\begin{proof}
Write $x-c=m+n$ with $m\in M$, $n\in N$.  The condition $Lx=t$
is equivalent to
\[
 Lm=t-Lc,
\]
so it forces $m=z_t$.  Orthogonality and Pythagoras give
\[
 \|x-c\|^2=\|z_t\|^2+\|n\|^2.
\]
This proves \eqref{eq:fiber}, as well as the descriptions of $I$ and $J$.

Assume first that $N\ne\{0\}$.  If $a,b\in M$, $\alpha,\beta\ge0$,
and $n\in N$, orthogonality gives
\begin{equation}\label{eq:pointdist}
 \dist(a+n,b+\beta\B_N)^2
 =\|a-b\|^2+(\|n\|-\beta)_+^2.
\end{equation}
Taking the supremum over $\|n\|\le\alpha$ shows
\[
 D_{\mathrm{asy}H}(a+\alpha\B_N,b+\beta\B_N)^2
 =\|a-b\|^2+(\alpha-\beta)_+^2.
\]
Interchanging the sets and taking the maximum proves \eqref{eq:exactH}.
When $N=\{0\}$, formula \eqref{eq:trivialH} is immediate.

Let $d=\|z_s-z_t\|$.  If $\|z_s\|,\|z_t\|\le b<r$, then
\begin{align*}
 |\rho_s-\rho_t|
 &=\frac{|\|z_t\|^2-\|z_s\|^2|}{\rho_s+\rho_t}\\
 &\le \frac{b}{\sqrt{r^2-b^2}}d.
\end{align*}
Combining this with \eqref{eq:exactH}, or just using
\eqref{eq:trivialH}, yields
\[
 D_H(A_s,A_t)\le
 \sqrt{1+\frac{b^2}{r^2-b^2}}\,d
 =\frac r{\sqrt{r^2-b^2}}d
 \le\frac{r\|G\|}{\sqrt{r^2-b^2}}\|s-t\|,
\]
which is \eqref{eq:localLip}.

For the boundary estimate, the elementary inequality
$(\sqrt u-\sqrt v)^2\le|u-v|$ gives
\[
 |\rho_s-\rho_t|^2
 \le |\|z_s\|^2-\|z_t\|^2|
 \le 2r d.
\]
Also $d\le2r$, hence $d^2\le2rd$.  Therefore
\[
 D_H(A_s,A_t)^2\le d^2+2rd\le4rd
 \le4r\|G\|\|s-t\|,
\]
with the singleton case even easier.  This proves \eqref{eq:holder} and all
continuity assertions.
\end{proof}

\section{Optimal-value stability without ambient compactness}

For $f:C\to\R$, write
\[
 \phi_-(t):=\inf_{x\in A_t}f(x),\qquad
 \phi_+(t):=\sup_{x\in A_t}f(x).
\]

\begin{theorem}[Infinite-dimensional optimal-value extension]\label{thm:values}
In the setting of Theorem~\ref{thm:geometry}:
\begin{enumerate}
\item If $f:C\to\R$ is uniformly continuous, then $f$ is bounded on
  $C$, and both $\phi_-$ and $\phi_+$ are finite and uniformly continuous
  on $J$.  If $f$ is
  $\Lambda$-Lipschitz, then
  \begin{equation}\label{eq:valueLip}
  |\phi_\pm(s)-\phi_\pm(t)|\le\Lambda D_H(A_s,A_t),
  \end{equation}
  so \eqref{eq:localLip} and \eqref{eq:holder} give explicit moduli.
\item If $N=\ker L$ is finite-dimensional and $f:C\to\R$ is merely
  continuous, then $\phi_-$ and $\phi_+$ are finite, attained, and
  continuous on $J$.
\end{enumerate}
Thus both compactness of $C$ and finite dimensionality of the ambient space
can be removed.  In particular, part 1 allows an arbitrary, possibly
infinite-dimensional kernel, while part 2 retains the source's mere
continuity hypothesis whenever the kernel is finite-dimensional.
\end{theorem}

\begin{proof}
First suppose that $f$ is uniformly continuous.  It is bounded on $C$:
choose $\delta>0$ such that $\|x-y\|<\delta$ implies
$|f(x)-f(y)|<1$, fix $x_0\in C$, and join $x_0$ to any $x\in C$
by a line segment split into more than $2r/\delta$ equal pieces.  The
triangle inequality bounds $|f(x)-f(x_0)|$ uniformly in $x$.

Now let $s_j\to s$ in $J$.  Theorem~\ref{thm:geometry} gives
$D_H(A_{s_j},A_s)\to0$.  Uniform continuity then implies
\[
 \inf_{A_{s_j}}f\longrightarrow\inf_{A_s}f,
 \qquad
 \sup_{A_{s_j}}f\longrightarrow\sup_{A_s}f.
\]
Indeed, approximate each point of either set by a point of the other and use
the uniform modulus; no minimizer is required.  Estimate \eqref{eq:holder}
makes this argument uniform on all of $J$.  For a $\Lambda$-Lipschitz
function it gives the standard estimate \eqref{eq:valueLip} directly, with
the sharper interior bound \eqref{eq:localLip}.

Suppose next that $N$ is finite-dimensional and $f$ is continuous.  Each
$A_t$ is compact by \eqref{eq:fiber}, so its two extrema are finite and
attained.  Given a sequence $s_j\to s$, the pairs
$(z_{s_j},\rho_{s_j})$ converge to $(z_s,\rho_s)$.  The set
\[
 P:=\{(z_s,\rho_s)\}\cup\{(z_{s_j},\rho_{s_j}):j\ge1\}
\]
is compact.  Therefore
\[
 K:=\bigcup_{j\ge1}A_{s_j}\cup A_s
 =\{c+z+\rho n:(z,\rho)\in P,\ n\in\B_N\}
\]
is the continuous image of the compact set $P\times\B_N$, and hence is
compact.  The restriction $f|_K$ is uniformly
continuous, and the first argument applied to $A_{s_j},A_s\subset K$
proves sequential continuity.  Since $J$ is metric, this is continuity.
\end{proof}

\section{Sharpness: continuity alone can fail}

The finite-dimensional-kernel condition in the second part cannot be removed
without some replacement such as uniform continuity.

\begin{proposition}[Escaping continuous wells]\label{prop:counterexample}
Let
\[
 H=\R\oplus\ell^2,\qquad Y=\R,\qquad L(a,x)=a,
 \qquad C=\overline B_H(0,1).
\]
Then $I=(-1,1)$ and
\[
 A_t=\{(t,x):\|x\|\le\sqrt{1-t^2}\}.
\]
There is a bounded continuous $f:C\to[-1,0]$ such that
\[
 \inf_{A_0}f=0,
 \qquad
 \inf_{A_{t_n}}f=-1
\]
for a sequence $t_n\downarrow0$.  Hence the optimal value function is
discontinuous at an interior parameter even though
$D_H(A_{t_n},A_0)\to0$.  Replacing $f$ by $-f$ gives the same failure
for supremal values.
\end{proposition}

\begin{proof}
Let $(e_n)$ be the standard orthonormal basis of $\ell^2$, and put
\[
 t_n:=2^{-(n+2)},\qquad
 \rho_n:=\sqrt{1-t_n^2},\qquad
 y_n:=(t_n,\rho_ne_n)\in A_{t_n},\qquad
 \delta_n:=t_n/3.
\]
The nearest point of $A_0=\{0\}\oplus\B_{\ell^2}$ to $y_n$ is
$(0,\rho_ne_n)$, so
\begin{equation}\label{eq:distA0}
 \dist(y_n,A_0)=t_n>\delta_n.
\end{equation}
For $m\ne n$, orthogonality gives
\[
 \|y_n-y_m\|^2=(t_n-t_m)^2+\rho_n^2+\rho_m^2.
\]
Since $t_n\le1/8$, the centers are uniformly separated, whereas
$\delta_n\le1/24$.  Thus the relative balls
$B(y_n,\delta_n)\cap C$ are pairwise separated and form a locally finite
family.

Define
\begin{equation}\label{eq:bump}
 f(x):=
 \begin{cases}
 -\left(1-\|x-y_n\|/\delta_n\right),
   &\text{if }\|x-y_n\|<\delta_n\text{ for some }n,\\
 0,&\text{otherwise.}
 \end{cases}
\end{equation}
The supports are disjoint and locally finite, and the displayed expression
vanishes on each support boundary; hence $f$ is continuous on $C$.
Equation \eqref{eq:distA0} implies $f=0$ on $A_0$, while
$f(y_n)=-1$.  Because $f\ge-1$, the asserted infima follow.  Notice also
that $x_n=(0,\rho_ne_n)\in A_0$ satisfies
$\|x_n-y_n\|=t_n\to0$ but
$|f(x_n)-f(y_n)|=1$; this pinpoints the failure of uniform continuity.
\end{proof}

\section{Verification, scope, and novelty}

The included checker audits the exact Hausdorff formula in finite-dimensional
proxies, both quantitative bounds, and the separation and distance claims in
Proposition~\ref{prop:counterexample}.  Run
\begin{verbatim}
conda run --no-capture-output -n sandbox python \
  code/check_hilbert_ball_fibers.py
\end{verbatim}
from the packet directory.  The computation is only a sanity check; all
arguments above are exact.

The source asks for removal ``at least in some interesting cases.''  The
closed-ball class gives a complete affirmative theorem: arbitrary Hilbert
dimension, closed-range operators, exact fiber distance, and sharp objective
hypotheses.  This is not a theorem for all closed bounded convex sets or all
Banach spaces.  The counterexample rules out the naive statement with an
arbitrary continuous objective and an infinite-dimensional kernel, but does
not rule out other geometric or coercive hypotheses.

Before promotion, the run indexes were searched for arXiv:1902.02363, the
source title, Hilbert balls, affine fibers, Hausdorff continuity, and
finite-dimensional kernels.  A bounded arXiv/web search used the exact source
sentence and close phrases involving affine sections of Hilbert balls and
noncompact parametric optimization.  It found the source paper and general
noncompact Berge-theorem extensions such as \cite{FKV,FKK}, but no exact
closed-ball fiber formula or sharpness theorem matching
Theorems~\ref{thm:geometry}--\ref{thm:values}.  This records search bounds; it
is not an exhaustive novelty certification.

Human review is recommended.  The highest-value checks are the exact
asymmetric Hausdorff calculation \eqref{eq:pointdist}, the compact-union
argument for finite-dimensional kernels, and local finiteness in the escaping
well construction.

\begin{thebibliography}{9}
\bibitem{RRD}
D. Reem, S. Reich, and A. De Pierro,
\emph{Stability of the optimal values under small perturbations of the
constraint set}, Pure Appl. Funct. Anal. 5 (2020), no. 3, 705--731;
arXiv:1902.02363v4.
\url{https://arxiv.org/abs/1902.02363}.

\bibitem{Conway}
J. B. Conway, \emph{A Course in Functional Analysis}, second edition,
Graduate Texts in Mathematics 96, Springer, 1990.

\bibitem{FKV}
E. A. Feinberg, P. O. Kasyanov, and M. Voorneveld,
\emph{Berge's maximum theorem for noncompact image sets},
J. Math. Anal. Appl. 413 (2014), 1040--1046; arXiv:1309.7708.

\bibitem{FKK}
E. A. Feinberg, P. O. Kasyanov, and D. N. Kraemer,
\emph{Continuity of parametric optima for possibly discontinuous functions
and noncompact decision sets}, arXiv:2109.06299, 2021.
\end{thebibliography}

\end{document}
